\section{Matrix types}

\subsection{By shape}

\index{matrices!row vector}
A \textit{row vector} is a matrix with one row.
\begin{equation*}
  \begin{bmatrix}
    1 & 2
  \end{bmatrix}
\end{equation*}

\index{matrices!column vector}
A \textit{column vector} is a matrix with one column.
\begin{equation*}
  \begin{bmatrix}
    1 \\
    2
  \end{bmatrix}
\end{equation*}

\index{matrices!square}
A \textit{square matrix} has the same number of rows and columns.
\begin{equation*}
  \begin{bmatrix}
    1 & 2 \\
    3 & 4
  \end{bmatrix}
\end{equation*}

\index{matrices!rectangular}
A \textit{rectangular matrix} has a different number of rows and columns.
\begin{equation*}
  \begin{bmatrix}
    1 & 2 & 3 \\
    4 & 5 & 6
  \end{bmatrix}
\end{equation*}

\subsection{By contents}

\index{matrices!zero}
A \textit{zero matrix} is a matrix filled with zeroes, denoted by $\mat{0}$.
\begin{equation*}
  \mat{0}_{2 \times 3} =
  \begin{bmatrix}
    0 & 0 & 0 \\
    0 & 0 & 0
  \end{bmatrix}
\end{equation*}

\index{matrices!one}

A \textit{one vector} is a column vector filled with ones, denoted by $\mat{e}$.
\begin{equation*}
  \mat{e} =
  \begin{bmatrix}
    1 \\
    1
  \end{bmatrix}
\end{equation*}

A \textit{one matrix} is a matrix filled with ones, denoted by $\mat{1}$.
\begin{equation*}
  \mat{1}_{2 \times 3} =
  \begin{bmatrix}
    1 & 1 & 1 \\
    1 & 1 & 1
  \end{bmatrix}
\end{equation*}

\index{matrices!identity}
An \textit{identity matrix} is a diagonal matrix of ones, denoted by $\mat{I}$.
\begin{equation*}
  \mat{I}_{3 \times 3} =
  \begin{bmatrix}
    1 & 0 & 0 \\
    0 & 1 & 0 \\
    0 & 0 & 1
  \end{bmatrix}
  \quad
  \mat{I}_{2 \times 3} =
  \begin{bmatrix}
    1 & 0 & 0 \\
    0 & 1 & 0
  \end{bmatrix}
\end{equation*}

\index{matrices!diagonal}
A matrix's \textit{diagonal} or \textit{main diagonal} is the list of entries
with the same row and column index.
\begin{equation*}
  \begin{bmatrix}
    {\color{red}1} & 0 & 0 \\
    0 & {\color{red}2} & 0 \\
    0 & 0 & {\color{red}3}
  \end{bmatrix}
\end{equation*}

\index{matrices!antidiagonal}
A matrix's \textit{antidiagonal} is the list of entries from the top-right to
the bottom-left perpendicular to the diagonal.
\begin{equation*}
  \begin{bmatrix}
    0 & 0 & {\color{red}1} \\
    0 & {\color{red}2} & 0 \\
    {\color{red}3} & 0 & 0
  \end{bmatrix}
\end{equation*}

A \textit{diagonal matrix} has nonzeroes restricted to the diagonal.
\begin{equation*}
  \begin{bmatrix}
    1 & 0 & 0 \\
    0 & 2 & 0 \\
    0 & 0 & 3
  \end{bmatrix}
\end{equation*}

\index{matrices!block diagonal}
A \textit{block diagonal matrix} has matrices along its diagonal and zeroes
elsewhere.
\begin{equation*}
  \begin{bmatrix}
    1 & 2 & 0 & 0 & 0 \\
    3 & 4 & 0 & 0 & 0 \\
    0 & 0 & 1 & 2 & 3 \\
    0 & 0 & 4 & 5 & 6 \\
    0 & 0 & 7 & 8 & 9
  \end{bmatrix}
\end{equation*}

\index{matrices!superdiagonal}
A matrix's \textit{superdiagonal} is the set of entries above and to the right
of the main diagonal.
\begin{equation*}
  \begin{bmatrix}
    0 & {\color{red}1} & 0 \\
    0 & 0 & {\color{red}2} \\
    0 & 0 & 0
  \end{bmatrix}
\end{equation*}

\index{matrices!subdiagonal}
A matrix's \textit{subdiagonal} is the set of entries below and to the left of
the main diagonal.
\begin{equation*}
  \begin{bmatrix}
    0 & 0 & 0 \\
    {\color{red}1} & 0 & 0 \\
    0 & {\color{red}2} & 0
  \end{bmatrix}
\end{equation*}

\index{matrices!banded}
\index{matrices!tridiagonal}
A \textit{banded matrix} has nonzero elements restricted to a diagonal band. A
\textit{tridiagonal matrix} has nonzero elements restricted to the main
diagonal, superdiagonal, and subdiagonal.
\begin{equation*}
  \begin{bmatrix}
    2 & 1 & 0 & 0 \\
    3 & 2 & 1 & 0 \\
    0 & 3 & 2 & 1 \\
    0 & 0 & 3 & 2
  \end{bmatrix}
\end{equation*}

\subsection{By structure}

\index{matrices!symmetric}
A \textit{symmetric matrix} is equal to its transpose.
\begin{equation*}
  \begin{bmatrix}
    1 & 2 \\
    2 & 3
  \end{bmatrix}
\end{equation*}

\index{matrices!orthogonal}
An \textit{orthogonal matrix} $\mat{Q}$ has the following properties:
\begin{itemize}
  \item $\mat{Q}\T\mat{Q} = \mat{Q}\mat{Q}\T = \mat{I}$
  \item $\mat{Q}\T = \mat{Q}^{-1}$
  \item $\det(\mat{Q}) = 1 \text{ or } -1$
\end{itemize}

\index{matrices!special orthogonal}
A \textit{special orthogonal matrix} is an orthogonal matrix with a determinant
of $1$.

\subsection{Matrix definiteness}

Tables \ref{tab:definiteness_eigenvalues} and \ref{tab:definiteness_definition}
describe the different kinds of matrix definiteness.
\begin{booktable}
  \begin{tabular}{|ccc|}
    \hline
    \rowcolor{headingbg}
    \textbf{Type} & \textbf{Eigenvalues} & \textbf{Notation} \\
    \hline
    Negative definite &
      All $\lambda < 0$ &
      $\mat{M} < \mat{0}$ \\
    Negative semidefinite &
      All $\lambda \leq 0$ &
      $\mat{M} \leq \mat{0}$ \\
    Positive semidefinite &
      All $\lambda \geq 0$ &
      $\mat{M} \geq \mat{0}$ \\
    Positive definite &
      All $\lambda > 0$ &
      $\mat{M} > \mat{0}$ \\
    Indefinite &
      Positive and negative &
      N/A \\
    \hline
  \end{tabular}
  \caption{Eigenvalue distribution and notation for each type of matrix
    definiteness. Let $\mat{M}$ be a matrix.}
  \label{tab:definiteness_eigenvalues}
\end{booktable}
\begin{booktable}
  \begin{tabular}{|cc|}
    \hline
    \rowcolor{headingbg}
    \textbf{Type} & \textbf{Definition} \\
    \hline
    Negative definite &
      $\mat{x}\T\mat{M}\mat{x} < 0$ for all $\mat{x}$ \\
    Negative semidefinite &
      $\mat{x}\T\mat{M}\mat{x} \leq 0$ for all $\mat{x}$ \\
    Positive semidefinite &
      $\mat{x}\T\mat{M}\mat{x} \geq 0$ for all $\mat{x}$ \\
    Positive definite &
      $\mat{x}\T\mat{M}\mat{x} > 0$ for all $\mat{x}$ \\
    Indefinite &
      $\exists \mat{x}, \mat{y} \ni \mat{x}\T\mat{M}\mat{x} < 0 < \mat{y}\T\mat{M}\mat{y}$ \\
    \hline
  \end{tabular}
  \caption{Rigorous definition for each type of matrix definiteness. Let
    $\mat{M}$ be a matrix and let $\mat{x}$ and $\mat{y}$ be nonzero vectors.}
  \label{tab:definiteness_definition}
\end{booktable}
